\chapter{Intelligente Trainingsdaten-Selektion fuer die Annotation historischer Karten}
\label{ch:training_data_selection}

\section{Motivation und Zielsetzung}
\label{sec:motivation_training_data}

Die Annotation historischer Karten ist teuer und zeitaufwaendig. Fuer die
vorliegende Masterarbeit ist deshalb nicht die zufaellige Auswahl moeglichst
vieler Kartenkacheln entscheidend, sondern die Konstruktion eines kleinen,
wissenschaftlich vertretbaren und reproduzierbar eingefrorenen Datensatzes.

Das Ziel des Dataselector-Workflows ist ein \emph{frozen dataset}, das

\begin{itemize}
    \item die primaere wissenschaftliche Stichprobe fuer die Arbeit definiert,
    \item die Auswahlentscheidung transparent dokumentiert,
    \item Fallbeispiele separat behandelt und
    \item fuer nachgelagerte Annotation und Modellvergleiche stabil bleibt.
\end{itemize}

\section{Kandidatenpool und Selektionsprinzip}
\label{sec:algorithmic_selection}

Der zugrunde liegende Kandidatenpool umfasst 673 georeferenzierte Kacheln der
\emph{Karte des Deutschen Reiches 1:100.000}. Die Selektion folgt einem
mehrkriteriellen Diversitaetsansatz:

\begin{equation}
\text{Score} = \alpha \cdot D_{\text{visual}} + \beta \cdot D_{\text{spatial}} + \gamma \cdot D_{\text{temporal}}
\end{equation}

Dabei stehen

\begin{itemize}
    \item $D_{\text{visual}}$ fuer Diversitaet im Merkmalsraum der Bildfeatures,
    \item $D_{\text{spatial}}$ fuer raeumliche Streuung und
    \item $D_{\text{temporal}}$ fuer zeitliche Streuung.
\end{itemize}

Die Gewichte $\alpha, \beta, \gamma$ sowie die Zielgroesse $n_{\text{samples}}$
werden im Parameter-Resolution-Schritt bestimmt und in einem validierten
Snapshot festgehalten. Fuer den aktuellen Thesis-v2-Vertrag gilt:
die Core-Auswahl wird aus genau diesen Snapshot-Parametern materialisiert
(\texttt{selection.selection\_authority = snapshot\_primary}).
Exploration und Autoscale verwenden dafuer dasselbe normalisierte Objective
(\texttt{selection.objective\_authority = unified\_normalized}).

\section{Aktueller Workflow des Thesis-Freezes}
\label{sec:current_workflow}

Der kanonische Workflow laeuft ueber
\texttt{python -m dataselector thesis-orchestrate}. Er trennt vier Ebenen:

\begin{enumerate}
    \item \textbf{Parameter Resolution}: Autoscale/Optuna bestimmen den
          Parameterkontext und schreiben einen validierten Snapshot.
    \item \textbf{Snapshot-primary Auswahl}: der Core-Datensatz wird direkt
          aus den Snapshot-Parametern neu materialisiert.
    \item \textbf{Core+Case-Vertrag}: primaere Core-Stichprobe und zusaetzliche
          Case-Tiles werden getrennt dokumentiert.
    \item \textbf{Patch-Handoff}: fuer die Annotation wird aus dem gefrorenen
          Tile-Satz ein Patch-Handoff erzeugt.
\end{enumerate}

Historische Runs koennen weiterhin einen Legacy-Modus verwenden
(\texttt{materialized\_csv\_primary}), in dem \texttt{tuning\_weights}-CSV-Dateien
als Auswahlquelle dienen. Dieser Modus ist fuer Rueckwaertskompatibilitaet
gedacht und nicht der Thesis-v2-Default.

\section{Eingefrorener Datensatz der Masterarbeit}
\label{sec:frozen_dataset}

Die Datensatzgroesse wird im Thesis-v2-Flow nicht als fixe Zahl im Text
angenommen, sondern direkt aus den Run-Artefakten gelesen:

\begin{itemize}
    \item \textbf{Core-Selektion}: \verb|selection_contract.json -> core_count|
    \item \textbf{Case-Tiles}: \verb|selection_contract.json -> case_count_attached|
    \item \textbf{Operative Gesamtmenge}: \verb|selection_contract.json -> final_count|
    \item \textbf{Patch-Handoff}: \verb|handoff_manifest.json -> selected_patch_count|
\end{itemize}

Die wissenschaftliche Interpretation ist dabei klar getrennt:

\begin{itemize}
    \item Primaere methodische Aussagen und Kernauswertungen beziehen sich auf
          die \textbf{Core-Selektion}.
    \item Hamburg wird als \textbf{Case-only}-Beispiel separat gefuehrt und
          erst danach an die operative Gesamtmenge angehaengt.
    \item Der Patch-Handoff ist ein abgeleitetes Annotation-Artefakt und aendert
          die Definition der zugrunde liegenden Tile-Stichprobe nicht.
\end{itemize}

\section{Reproduzierbarkeit und Provenienz}
\label{sec:reproducibility}

\textbf{Dataset authority:} \verb|selection_core.csv|,
\verb|selection_final_with_cases.csv| und
\verb|selection_contract.json| definieren den eingefrorenen
Thesis-Datensatz.

\textbf{Parameter authority:} der validierte Snapshot und die
Parameter-Resolution-Artefakte definieren den aufgeloesten Parameterkontext.

Fuer die Arbeit ist diese Unterscheidung zentral. Im kanonischen Thesis-v2-Run
wird erwartet, dass Snapshot-Parameter und materialisierte Auswahl konsistent
sind (\texttt{Selection Reconciliation = aligned}). Fuer historische Legacy-Runs
gelten folgende Regeln:

\begin{itemize}
    \item Datensatz-Claims binden an die materialisierten Auswahl-Artefakte.
    \item Parameter-Claims binden an Snapshot-/Resolver-Artefakte.
    \item Abweichungen muessen im Report explizit als
          \texttt{documented\_difference} markiert werden.
\end{itemize}

Die Provenienz des aktuellen Freezes wird konkret ueber folgende Artefakte
dokumentiert:

\begin{itemize}
    \item \verb|selection_core.csv|
    \item \verb|selection_case.csv|
    \item \verb|selection_final_with_cases.csv|
    \item \verb|selection_contract.json|
    \item \verb|final_config_<timestamp>.yaml|
    \item \verb|parameter_resolution/optuna_autoscale_best_latest.json|
\end{itemize}

\section{Wissenschaftliche Vertretbarkeit}
\label{sec:scientific_justification}

Der aktuelle Freeze ist wissenschaftlich vertretbar, weil

\begin{itemize}
    \item die primaere Stichprobe explizit als Core-Datensatz definiert ist,
    \item Fallbeispiele die Core-Auswertung nicht konfundieren,
    \item Unsicherheitsquantifizierung ueber
          \texttt{bootstrap\_candidates} erfolgt und
    \item die finale Stichprobe als archiviertes Run-Artefakt reproduzierbar
          vorliegt.
\end{itemize}

Die entscheidende methodische Aussage lautet daher nicht
\glqq diese Gewichte erzeugten exakt diese Tiles\grqq, sondern
\glqq diese Masterarbeit wertet einen konkret eingefrorenen, dokumentierten und
reproduzierbaren Datensatz aus\grqq.

\section{Zusammenfassung}
\label{sec:summary_training_data}

Die Masterarbeit verwendet einen methodisch klar definierten, aus Run-Artefakten
abgeleiteten Frozen-Datensatz. Die wissenschaftliche Kerninterpretation basiert
auf der Core-Selektion; optionale Case-Tiles werden transparent getrennt
dokumentiert. Durch die explizite Trennung von Dataset authority und Parameter
authority sowie den Snapshot-primary-v2-Vertrag bleibt die Auswahl kausal
rueckfuehrbar und reproduzierbar.
